% =============================================================
% Annexe E — Protocoles et preuves de validation
% Alimentée depuis Ch.6 (§6.1, §6.2).
% Absorbe, depuis l'ancienne annexe H démontée :
%   - H.1 « preuves complémentaires disponibles » (2 mesures datées),
%   - H.3 « règles de production » resserrées en note de méthode.
% Aucun appel \capture : aucune image n'existe dans latex/figures/.
% Condensation vague 3 : le tableau des conditions de mesure, le tableau des
% mesures complémentaires et la section de synthèse finale sont supprimés
% (redites du ch.6) ; l'analyse de T1, T8 et T11 reste au corps, l'annexe n'en
% garde que le protocole, le critère et la portée de la preuve.
% =============================================================
\chapter{Protocoles et preuves de validation}
\label{ann:protocoles-tests}

Cette annexe est la base de reproduction de la campagne de vingt protocoles dont le
\cref{ch:validation} publie les résultats : un tiers disposant du même environnement peut
rejouer chacun d'eux, son critère d'acceptation étant formulé avant exécution et publié ici. Elle
donne les conditions de mesure, le modèle de fiche appliqué à chaque protocole, et les vingt
protocoles avec leur critère, leur résultat et sa date. Deux relevés de recette les complètent sans
en fonder aucun --- antériorité du dernier rattachement sémantique sur l'attaque du 19/08/2026,
durées de retour arrière applicatif du 07/08/2026 --- et sont publiés avec leur date en
\cref{sec:eval-enrichissement,sec:performance}.

\medskip
\noindent\textbf{Conditions de mesure et modèle de fiche.} La campagne est exécutée sur l'environnement de recette, de topologie identique à la cible
(\cref{sec:env-travail}) ; les preuves ponctuelles sont datées jusqu'au 25/08/2026, les preuves
automatisées sont rejouées à chaque livraison et n'ont pas de date unique de relevé. Deux
paramètres ne sont pas épinglés, et le sont dits : la version du référentiel de configuration
d'infrastructure n'est pas consignée, et la matrice de techniques d'attaque non plus ---
l'instantané employé est identifié par l'empreinte publiée en \cref{sec:eval-detection}. Aucun taux
de conformité ni de couverture n'en est tiré.

Chaque protocole suit le même modèle : un identifiant \textbf{Ti} renvoyant à l'exigence
\textbf{EXi} et à la menace qu'il vérifie ; un \textbf{protocole} décrivant l'action tentée ; un
\textbf{critère d'acceptation} formulé avant l'exécution et non révisé après coup ; un
\textbf{résultat} selon la convention à quatre valeurs de \cref{sec:strategie} ; et la
\textbf{preuve} qui le fonde, avec sa date ou son mode d'exécution. Les six protocoles rattachés au
plan de validation continue portaient en outre un \textbf{attendu} --- ce que la preuve devrait
rendre visible, énoncé \emph{avant} l'exécution --- reproduit ci-dessous sans révision, à côté du
relevé qui l'a confronté : l'écart est nul pour T17 et T20, partiel pour T12, T15 et T19,
défavorable pour T14.

Une règle de méthode gouverne enfin la production de ces preuves : \textbf{la sortie brute
d'exécution est conservée, horodatée dans le cadre même où elle est produite}, et c'est elle qui
fait preuve --- une date portée par une légende est une affirmation de l'auteur, une date lisible
dans la sortie est une propriété de la preuve. Les valeurs sensibles qu'elle porte y sont masquées
par rectangle opaque et jamais par floutage, un flou étant réversible par traitement d'image.

\medskip
\noindent\textbf{Protocoles détaillés.}

\textbf{T1 --- Filtrage des attaques applicatives} (EX1, SO1). Envoyer, depuis l'extérieur, des
charges d'attaque correspondant aux cinq familles de règles activées au périmètre.
\emph{Critère} : toutes les requêtes sont rejetées avant d'atteindre la charge de travail, verdict
journalisé. \emph{Résultat} : \textbf{non conforme}. \emph{Preuve} : le 23/08/2026, les charges
d'injection SQL, de script inter-sites et d'inclusion de fichier local ont été refusées en~403,
mais une forme brute de traversée de chemin a reçu une redirection~302 --- le quantificateur du
critère n'est pas tenu ; mécanisme et correctif en \cref{sec:non-conformite-t1}.

\textbf{T2 --- Limitation des tentatives d'authentification} (EX2, SO2/SO3). Générer des échecs
d'authentification répétés depuis une adresse unique jusqu'au franchissement du seuil.
\emph{Critère} : blocage temporaire déclenché, trace présente dans l'entrepôt. \emph{Résultat} :
conforme. \emph{Preuve} : le 19/08/2026, onze requêtes ont produit dix réponses~422 puis une
réponse~429.

\textbf{T3 --- Comportement sous charge et traçabilité du balayage} (EX3, SO4/SO5). Générer une
charge supérieure au seuil global puis un balayage de la surface exposée. \emph{Critère} : le
service reste disponible, les deux activités apparaissent dans les journaux. \emph{Résultat} :
partiellement conforme. \emph{Preuve} : le balayage est tracé, les journaux du répartiteur étant
collectés sans filtre de statut ni de sévérité. \emph{Réserve} : l'essai de charge au-delà du seuil
global n'est pas documenté.

\textbf{T4 --- Inaccessibilité des données depuis Internet} (EX4, SO6). Depuis un hôte externe,
tenter une connexion directe à la base de données ; vérifier par inventaire qu'aucune adresse
publique n'est attribuée. \emph{Critère} : aucune adresse publique, connexion en échec.
\emph{Résultat} : partiellement conforme. \emph{Preuve} : l'inventaire de l'instance montre
l'attribution d'adresse publique désactivée, un raccordement par accès privé aux services et un
mode de transport strictement chiffré. \emph{Réserve} : l'échec de connexion est \emph{déduit} de
l'absence de route, non relevé depuis un hôte externe --- c'est la déduction, non la mesure, qui
borne ce résultat.

\textbf{T5 --- Cloisonnement des identités applicatives} (EX5, SO7). Avec l'identité d'une
application hébergée, tenter une lecture de l'entrepôt de supervision. \emph{Critère} : refus
d'autorisation. \emph{Résultat} : conforme. \emph{Preuve} : deux rôles exactement attribués et usurpation refusée
en direct le 23/08/2026 ; l'identité du composant d'encodage est dépourvue de tout rôle sur
l'entrepôt, vérifié le 19/08/2026. Un test automatisé rejoue à chaque livraison la lecture
interdite du secret de base applicative sous l'identité d'enrichissement et exige le refus.

\textbf{T6 --- Authentification des appels entre services} (EX6, SO8). Appeler le service
d'encodage sans jeton d'identité. \emph{Critère} : refus --- être « à côté » ne doit rien
accorder. \emph{Résultat} : partiellement conforme. \emph{Preuve} : un test automatisé, en porte bloquante à
chaque livraison, émet l'appel sans jeton et vérifie le refus ; l'entrée du service est restreinte
aux appels internes et le droit d'invocation à la seule identité d'enrichissement.
\emph{Réserve} : l'appel du test vient de l'extérieur du réseau privé, et le refus atteste donc le
contrôle de périmètre ; l'appel depuis l'intérieur sans jeton, formulation exacte du critère, reste
porté par la configuration.

\textbf{T7 --- Absence de clés de longue durée} (EX7, SO9). Inventorier les clés existantes pour
toutes les identités de service ; tenter d'en créer une. \emph{Critère} : inventaire vide et
absence de clé exportée dans la chaîne. \emph{Résultat} : conforme. \emph{Preuve} : le critère est
un inventaire, et l'inventaire est vide --- recherche exhaustive sur le dépôt et l'infrastructure,
aucune ressource de clé de compte de service, aucun fichier d'identifiants, aucune variable de
justificatifs ; la fédération d'identité, contrainte à la fois sur le dépôt et sur la branche, rend
la clé inutile. \emph{Qualification} : la garantie est \emph{vérifiée} par inventaire reproductible
et non \emph{préventive}, l'interdiction structurelle de créer une clé relevant d'une politique
d'organisation hors du périmètre administratif disponible (écart É7). La distinction qualifie le
résultat sans l'annuler : le critère ne porte pas sur l'existence d'un contrôle préventif.

\textbf{T8 --- Blocage d'un artefact vulnérable} (EX8, SO10). Soumettre à la chaîne une image
contenant une vulnérabilité critique connue. \emph{Critère} : arrêt à l'étape d'analyse de
vulnérabilité \cite{slsa_framework}. \emph{Résultat} : conforme, sur un refus réellement survenu et
sans mise en scène. \emph{Preuve} et \emph{réserve} --- vulnérabilité de \texttt{node-tar} 6.2.1,
correctif retenu, ordre publication/analyse à réaligner --- en \cref{sec:preuve-t11}.

\textbf{T9 --- Registre de confiance unique} (EX9, SO11). Tenter un déploiement depuis une source
d'images externe, puis par référence mutable au lieu d'une empreinte. \emph{Critère} : les deux
tentatives échouent \cite{slsa_framework} --- une étiquette contenant une empreinte de commit ne
constitue pas une empreinte d'image immuable. \emph{Résultat} : partiellement conforme.
\emph{Preuve} : un registre unique, une résolution par empreinte d'image sur le chemin de
description en code, une étiquette par empreinte de commit sur le chemin de livraison --- immuable
en pratique, mutable en théorie. \emph{Réserve} : aucun contrôle n'oppose de refus, ni étiquettes
immuables ni admission d'image n'étant déclarées, et l'application hébergée est déployée par
étiquette flottante (écart É8). La propriété tient par construction, non par imposition : c'est la
borne exacte de ce résultat.

\textbf{T10 --- Blocage d'un secret versionné} (EX10, SO12). Soumettre un secret factice de format
reconnaissable ; vérifier séparément qu'aucun secret du système n'est dérivé d'un autre.
\emph{Critère} : arrêt à la première étape, indépendance confirmée par inventaire.
\emph{Résultat} : partiellement conforme. \emph{Preuve} : la porte de détection de secrets est
bloquante à chaque livraison comme à chaque demande de fusion, sur l'historique complet du dépôt,
une profondeur de clone complète lui étant explicitement accordée ; l'indépendance est établie par
inventaire, chaque secret étant engendré par un tirage aléatoire indépendant, stocké séparément et
injecté par référence, aucun mot de passe ne transitant par une variable de configuration.
\emph{Réserve} : aucun refus n'a été observé sur un secret factice, et le cas de la porte déclarée
bloquante mais inerte (\cref{subsec:incidents}) est la raison pour laquelle une porte non éprouvée
n'est pas comptée conforme.

\textbf{T11 --- Intégrité de la chaîne de preuve} (EX11, SO13/SO16) --- \emph{test le plus
discriminant de la campagne}. Avec l'identité de la tâche d'enrichissement, tenter une écriture
puis une suppression dans les tables de journaux bruts et de détections. \emph{Critère} : refus
sur ces tables, écriture restant possible sur la seule table d'enrichissement. \emph{Résultat} :
conforme. \emph{Preuve} : un test de bout en bout automatisé, rejoué à chaque livraison, tente
réellement l'insertion interdite sous l'identité usurpée et n'est vert que si elle est refusée ; la
vérification des droits table par table du 19/08/2026 couvre en outre la table des journaux bruts et
le cas de la suppression. Le mécanisme, l'histoire de ce test et sa \emph{réserve} --- tâche non
bloquante, abstention explicite lorsque l'identité courante n'a pas le droit d'usurpation --- sont
analysés en \cref{sec:preuve-t11}.

\textbf{T12 --- Détection d'une interruption de collecte} (EX12, SO14). Interrompre volontairement
la collecte au-delà du seuil d'alerte. \emph{Critère} : une alerte est émise. \emph{Attendu, écrit
avant exécution} : la notification effectivement reçue sur le canal déclaré, avec l'heure
d'interruption, l'heure de déclenchement et le délai qui les sépare. \emph{Résultat} :
partiellement conforme. \emph{Preuve} : le 22/08/2026, le générateur de trafic de recette a été arrêté à \hms{10:35:00} ;
la politique d'absence de signal s'est ouverte à \hms{11:10:44} et la notification a été reçue à
\hms{11:11:19}, soit 36~min~19~s après l'interruption --- 30~min de durée de condition, 300~s de
période d'alignement, puis évaluation et acheminement. \emph{Réserve} : exécutée d'abord sur
l'autre panne que recouvre le protocole, l'interruption de l'acheminement des journaux, la mesure
est défavorable --- l'acheminement coupé de \hms{09:12} à \hms{10:20} n'a ouvert aucune alerte,
1\,148 requêtes comptées par la métrique d'exécution n'ont produit aucune ligne dans l'entrepôt, et
le rétablissement n'a rien rejoué : la politique porte sur la production des journaux, jamais sur
leur acheminement ni sur la fraîcheur de la table qui les reçoit.

\textbf{T13 --- Attribution des actions privilégiées} (EX13, SO15). Exécuter une action
privilégiée identifiable, la retrouver dans les journaux d'audit avec son auteur. \emph{Critère} :
action retrouvée, auteur identifié, horodatage cohérent. \emph{Résultat} : partiellement conforme.
\emph{Preuve} : la journalisation d'audit des accès aux données est activée en lecture \emph{et} en
écriture sur les quatre services porteurs de données --- entrepôt, gestionnaire de clés,
gestionnaire de secrets, base relationnelle --- et le journal d'activité d'administration est
produit par construction, sans désactivation possible. \emph{Réserve} : ces journaux ne sont
acheminés vers aucune destination dédiée, demeurent dans un espace à rétention de trente jours
modifiable, et aucune règle ne les exploite. L'attribution est possible, elle n'est pas rendue
incontestable.

\textbf{T14 --- Sortie du composant d'inférence} (EX14, SO17). Depuis le composant d'encodage,
tenter une connexion sortante vers une destination externe. \emph{Critère} : échec de connexion.
\emph{Attendu, écrit avant exécution} : l'échec horodaté d'une connexion sortante émise depuis le
composant lui-même, et la trace de refus correspondante. \emph{Résultat} : \textbf{non conforme}.
\emph{Preuve} : le 25/08/2026, deux sondes ne différant que par leur chemin de sortie ont été
exécutées sous l'identité du composant. La sonde A, placée dans le sous-réseau plateforme avec
l'étiquette du composant et la sortie directe, a été refusée deux fois --- à \hms{10:04:14} vers
\texttt{203.0.113.10:443} et à \hms{10:04:26} vers \texttt{198.51.100.7:443} --- par la règle de
refus par défaut en sortie, chaque refus portant sa ligne de journal dans l'entrepôt et son code 28
côté client 10,3~s plus tard. La sonde B, portant la configuration réseau exacte de la révision
déployée --- sortie managée par la plateforme d'exécution, sans accès au réseau privé --- a reçu un
code~204 en 412~ms à \hms{10:12:38} ; la lecture de la révision à \hms{10:18} confirme l'absence de
bloc d'accès au réseau. Sur le chemin réellement emprunté par le composant, la connexion aboutit et
aucune règle de pare-feu ne s'applique. Le correctif, son coût et sa portée --- il vaut identiquement pour l'interface web de l'application
hébergée --- sont analysés en \cref{sec:non-conformite-t1}. \emph{Réserve} : le relevé porte sur la
capacité de sortie, non sur son usage --- la sonde B partage l'identité et le chemin réseau de la
révision déployée, non son binaire, et l'absence d'appel sortant dans l'image n'est établie que par
lecture du code, preuve de rang 4 qui ne peut fonder aucun verdict de conformité.

\textbf{T15 --- Comportement de la chaîne d'enrichissement sous charge} (EX15, SO18). Injecter un
volume de détections nettement supérieur au volume nominal. \emph{Critère} : la tâche se termine
sans perte, coût et durée mesurés. \emph{Attendu, écrit avant exécution} : le volume injecté, le
nombre de détections traitées, la durée de la tâche et l'égalité entre entrées et sorties, relevés
sur une même exécution. \emph{Résultat} : partiellement conforme. \emph{Preuve} : le 22/08/2026,
500 détections --- environ dix-sept fois le volume nominal --- ont été injectées à \hms{14:00:00} ;
sept exécutions successives, de \hms{14:15} à \hms{15:45}, ont lu 350 détections et écrit
350 lignes dans la table d'enrichissement, entrées égales aux sorties sur chaque exécution, aucune
en échec et aucune reprise. La durée maximale est de 96~s, dont 31~s de démarrage à froid du
service d'encodage, pour un délai d'expiration de 600~s ; la durée cumulée est de 519~s et le coût
marginal reste inférieur à 0,05~euro, dominé par le minimum forfaitaire de balayage par table et
non par le volume. \emph{Réserve} : le critère est tenu au niveau de la tâche, non de la chaîne ---
150 des 500 détections injectées sont sorties de la fenêtre glissante de deux heures avant d'avoir
été lues, sans erreur, sans code de retour non nul et sans trace. Le débit plafond mesuré est de 50 détections par quart d'heure, soit 200 par heure et une capacité
d'absorption de 400 détections en attente ; la politique « Enrichissement en retard » s'est ouverte
à \hms{14:35:47} et refermée à \hms{16:05:12}, non parce que le retard était résorbé mais parce que
les lignes en attente étaient sorties de la fenêtre --- l'alerte se referme à l'instant où la perte
devient définitive. Le trafic de recette ayant été suspendu pendant la mesure, ce débit est une
borne supérieure que le flux nominal, partageant le même lot de 50, abaisse.

\textbf{T16 --- Détection de dérive} (EX16, SO20). Modifier hors code un attribut réversible d'une
ressource déjà gérée dans l'état, puis exécuter une planification avec rafraîchissement.
\emph{Critère} : la planification signale la différence et les journaux d'audit identifient
l'auteur \cite{hashicorp_drift} ; une ressource nouvelle absente de l'état exige un contrôle
d'inventaire séparé. \emph{Résultat} : partiellement conforme. \emph{Preuve} : mécanisme sollicité en conditions réelles
lors de l'incident de chiffrement du 08/08/2026, planification sans différence obtenue le
19/08/2026, contrôle quotidien de dérive de l'isolation entre locataires doté de sa propre
politique d'alerte. \emph{Réserve} : aucune détection de dérive n'est planifiée sur l'infrastructure
décrite en code, la chaîne de provisionnement (F7) ne comportant ni planification ni application
automatisées, et une ressource créée hors état échapperait à la seule planification.

\textbf{T17 --- Protection de l'état de l'infrastructure} (EX17, SO21). Tenter un accès en lecture
à l'état avec une identité non autorisée. \emph{Critère} : refus. \emph{Attendu, écrit avant
exécution} : le refus d'autorisation horodaté, avec l'identité employée et le droit manquant nommé.
\emph{Résultat} : conforme. \emph{Preuve} : le 20/08/2026, quatre refus ont été opposés aux trois
identités de service susceptibles d'atteindre l'emplacement de l'état --- lecture refusée à
\hms{08:41:07} sous l'identité d'API, à \hms{08:42:15} sous celle de la tâche d'enrichissement et à
\hms{08:43:02} sous celle du tableau de bord, chacune en erreur~403 nommant le droit
\texttt{storage.objects.get} manquant, puis énumération du préfixe refusée à \hms{08:43:44} sur
\texttt{storage.objects.list}. Un contrôle positif à \hms{08:45:20}, sous l'identité de la chaîne
de livraison, énumère l'objet d'état : sans lui, quatre échecs ne distingueraient pas un refus d'un
chemin erroné. Preuve de rang 3, vérification en direct datée. \emph{Réserve} : l'emplacement de l'état est
provisionné hors description en code --- ni son versionnement, ni son mode d'accès uniforme, ni sa
politique de conservation ne sont déclarés, de sorte qu'une modification de ces propriétés ne serait
signalée par aucun contrôle de dérive ; le refus établit un état à un instant, non sa persistance,
d'où la cadence de rejeu et le déclencheur sur tout changement d'emplacement ou de droits inscrits
au plan. Le verrouillage concurrent est natif à l'emplacement retenu ; son versionnement n'a jamais
été éprouvé en reprise.

\textbf{T18 --- Revue avant modification de droits} (EX18, SO22). Soumettre une modification
élargissant les droits d'une identité. \emph{Critère} : blocage en attente d'approbation, aucune
application sans validation humaine. \emph{Résultat} : partiellement conforme. \emph{Preuve} : la revue de demande de fusion est en
place et l'infrastructure est appliquée manuellement par l'administrateur, ce qui interpose de fait
une validation humaine. \emph{Réserve} : cette validation n'est pas outillée --- aucun environnement
protégé n'est déclaré sur les six tâches de la chaîne de livraison, donc aucune approbation n'est
\emph{exigée} par un mécanisme. Le contrôle tient par la procédure, non par la porte.

\textbf{T19 --- Détection d'une consommation anormale} (EX19, SO23). Provoquer une consommation de
ressources nettement supérieure au profil habituel. \emph{Critère} : alerte émise.
\emph{Attendu, écrit avant exécution} : la notification reçue, le seuil franchi et le poste de
consommation qui l'a provoqué. \emph{Résultat} : partiellement conforme. \emph{Preuve} : le prérequis a été levé le 20/08/2026 par la déclaration, dans le code
d'infrastructure, d'un budget mensuel de 50,00~euros à seuils de 50\,\%, 90\,\% et 100\,\% de la
dépense réelle, acheminé vers les deux canaux déjà déclarés. Le 21/08/2026, le budget a été ramené
à 6,50~euros --- la dépense du mois s'élevant alors à 6,42~euros, soit 98,8\,\% --- puis une charge
de 2~h~30 a porté le générateur de trafic à vingt fois son débit et rejoué 150 fois chacune des
douze requêtes planifiées, soit 1\,800 exécutions et environ 41~Go balayés. La dépense de fin de
journée atteint 6,63~euros, 102,0\,\% du budget temporaire, pour 0,21~euro de consommation
additionnelle ; les notifications des seuils 90\,\% et 100\,\% ont été reçues à \hms{12:47} et
\hms{17:38}. Le critère « alerte émise » est tenu : deux notifications reçues sur le canal déclaré, pour deux
seuils distincts, avec leur horodatage. \emph{Réserve} : les deux autres éléments de l'attendu ne le
sont pas. Le délai entre franchissement et notification est un \emph{intervalle} compris entre
5~h~48 et 8~h~18 --- l'instant du franchissement n'étant connu qu'à l'heure près --- et il n'est
borné par aucun engagement ; la notification porte le nom du budget, le montant et le pourcentage,
jamais le poste de consommation, reconstitué après coup depuis la ventilation de coût. Le
dispositif est un signal de coût, non un signal de sécurité : un détournement de ressources
fonctionne plusieurs heures avant le départ de la première notification.

\textbf{T20 --- Restriction des sorties réseau} (EX20, SO24). Depuis une charge de travail
applicative, tenter une connexion sortante vers une destination non autorisée. \emph{Critère} :
échec de connexion et trace de refus exploitable. \emph{Attendu, écrit avant exécution} : l'échec
de la connexion sortante côté charge de travail, et la ligne de refus correspondante retrouvée dans
l'entrepôt de supervision, les deux horodatages se recoupant. \emph{Résultat} : conforme.
\emph{Preuve} : la condition de recevabilité était acquise --- journalisation des refus réseau
activée et vérifiée depuis le 19/08/2026, ce qui clôt l'écart É1. Le 24/08/2026, après la pose du
refus par défaut en sortie du 23/08 et avant le démontage de l'ancien chemin de traduction
d'adresses, une tâche déployée avec l'identité de service, l'étiquette, le sous-réseau et le mode
de sortie d'un locataire réel a émis trois tentatives, toutes refusées : à \hms{09:12:33} vers une
adresse externe non autorisée, par la règle de refus par défaut en sortie ; à \hms{09:13:49} vers
le sous-réseau de la plateforme, par la règle de refus croisé entre locataires ; à \hms{09:15:02}
vers l'adresse privée de l'instance de base sur le port 5432, par la règle nommée qui l'interdit.
Les trois lignes de refus ont été retrouvées dans l'entrepôt avec la règle, la source, la
destination et le port, interrogeables à \hms{09:18}, de 2~min~55~s à 4~min~12~s après l'événement,
chacune précédant l'échec client de 10,2 à 10,3~s --- le délai d'expiration demandé plus
l'acheminement. Un contrôle positif à \hms{09:16:41}, la même sonde portant l'étiquette de la tâche
d'enrichissement, atteint la destination autorisée par sa règle et reçoit un code~200 : sans lui,
trois échecs ne distingueraient pas un refus du pare-feu d'une sortie réseau en panne. Preuve de
rang 3, adossée à une configuration déclarée en code. \emph{Réserves} : elles bornent la portée, non la couverture du critère. L'étiquette se pose dans
la définition de la révision, de sorte que la frontière est opposable à une charge compromise mais
non à un opérateur capable de déployer ; seuls les refus sont journalisés, par décision de coût ;
l'exfiltration par les interfaces du fournisseur reste ouverte, l'autorisation qui rend l'entrepôt
joignable rendant joignable tout entrepôt du même fournisseur ; et deux charges ne sont pas
derrière ces règles --- le composant d'encodage et l'interface web de l'application hébergée ---
ce dont T14 est la conséquence directe.
